\subsection{Interface to the knowledge extraction pipeline}\label{subsec:doc-interface-to-knowledge-pipeline}


The knowledge extraction pipeline, described in more detail in Section~\ref{section:Knowledge_Extraction}, consumes a very narrow interface, consisting of a stream of sentence-segmented text, where each sentence is tagged with the database identifier of the text element it was taken from. The figures, tables, cross-references, and section hierarchy are also persisted by the document extraction pipeline, but are outside the scope of this thesis; the knowledge extraction pipeline can be extended to include visual language models, and process the table and figure crops as well. Therefore, the knowledge pipeline, with its current implementation, only reads text elements.

\pinktodo{Somewhere in the thesis, possibly in the Future Work section, talk about adding more context to the knowledge extraction pipeline, by incorporating the text extracted from tabls and figures, and also some text context from the rest of the section.}

Concretely, the MAP stage, described in Subsection~\ref{subsec:knowledge-map}, reads the text elements of a single paper, which are sentence-segmented and grouped into chunks. Each sentence reaches the voter as a citation-tagged line, in the following format: \texttt{[$S_\textit{i}$|PMCID|\textit{te\_id}]}, followed by the sentence text. \texttt{te\_id} is the primary key of the source element, and is the only structured context the voter receives: the sentence string and its source identifier, with no section title, position in section, or links to the figures or tables it refers to. The section context and the full paragraph remain retrievable through \texttt{te\_id}, but at the moment, they are not part of what the voters in the knowledge extraction pipeline sees.

Despite its narrowness, this interface is what makes document extraction more than just pre-processing. The \texttt{te\_id} on each sentence is the anchor for the provenance chain, as every MAP finding records a verbatim source quote and citation identifiers of the sentences they use. Then, these identifiers go through normalization, grouping, canonicalization, and resolution, all of which will be explained in Section~\ref{section:Knowledge_Extraction}. Thanks to the identifiers, every final scored rule can be traced back to the text element rows they were inferred from. 

Since every downstream claim resolves to extracted text elements, errors in text extraction, sentence segmentation, text filtering and paragraph assembly can propagate into the findings and rules generated later. Therefore, the quality of the knowledge extraction output depends on the quality of the document extraction output. Chapter~\ref{chapter:Discussion} discusses this dependency as part of the general validity analysis of the entire system.

